%
% frame_info_popup_stale_press_buggy.tex
%
% FIDELITY CONTROL for frame_info_popup.tex -- NOT the design.
%
% This is the specification with exactly one rule changed, so that a
% defect the design closes becomes reachable.
%
% The edit:
%   FrameRelease clears the press record only when it judged a click;
%   any other release leaves it in place.
%
% Must find:
%   probcli frame_info_popup_stale_press_buggy.tex -model_check -nodead \
%     -goal "i4 : bad" -p DEFAULT_SETSIZE 2 -p MAX_OPERATIONS 1000
%

\documentclass[a4paper,10pt,fleqn]{article}
\usepackage[margin=1in]{geometry}
\usepackage{fuzz}
\usepackage[colorlinks=true,linkcolor=blue,citecolor=blue,urlcolor=blue]{hyperref}

\begin{document}

\title{The lux Frame-Info Popup: Fidelity Control, (d), the Press Record Kept After a Release That Was Not a Click}
\author{Formal model of the glyph press / hold / release state machine}
\date{October 2026}
\maketitle

\paragraph{This file is a fidelity control.}
It is \texttt{frame\_info\_popup.tex} with one rule changed ((d), the Press Record Kept After a Release That Was Not a Click), so that
the defect the design closes is reachable.  Where the prose below describes
the intact rule, the edited schema governs.  See Section~\ref{sec:fidelity} of
the specification for the trace it must reproduce.

% ============================================================
\section{Introduction}
% ============================================================

Every frame's title bar carries a small ``i'' glyph.  A click on it opens a
popup of the frame's owning-connection facts; a second click closes it.  That
sentence hides a good deal of machinery.  The glyph is not an ImGui item: it is
painted on the foreground draw list and hit-tested against the raw mouse
position, so ImGui knows nothing of it.  A press on it is, to ImGui, a press on
the frame's title bar, and a press on a title bar starts a window move.  While
that move is live ImGui re-focuses the frame on every frame, and focusing a
frame closes every popup stacked over it --- so the code's documentation says,
and Section~\ref{sec:assume} records how much of that could be confirmed.  A
popup opened on the \emph{press} would then not survive the held button; a
popup opened on the \emph{release}, after the move has ended, does.  So the glyph
must judge a click, not a press.

A click is the release over a glyph of a press that began on \emph{that}
glyph and never strayed past the drag threshold.  Each clause of that
definition has been wrong in the code at some point, and each time the popup
did something the user had not asked for or nothing at all:
\begin{itemize}
  \item The popup was opened on the press, and ImGui's per-frame refocus shut
    it before the button came up: the popup was never visible for a held
    click.
  \item The drag test compared the \emph{current} offset from the press point
    with the threshold, so dragging away and returning before the release
    opened the popup although the frame had been moved.
  \item The toggle was decided on the release from the popup's \emph{current}
    state.  ImGui shuts a popup on a press outside it --- a press on the glyph
    that the popup does not cover is such a press --- so by the release the
    popup was already closed, and the second click reopened what it had just
    seen close.
  \item A single press record, kept across glyphs and cleared only by a
    click, survived a release that was not one; a later release over some
    glyph whose press had never begun there was misjudged against it.
\end{itemize}
Four defects of one family across successive review rounds is the recurrence
signal: the problem is not a missing test but a state space, and the state space
is small enough to enumerate.  This document is that enumeration.

The model is a finite state machine over a bounded set of frames.  It has the
code's own state --- the popup, and the one press record the controller holds
--- and, beside it, the \emph{physical gesture} the user is making: whether the
button is down, where it went down, whether the pointer has strayed.  The
physical gesture is the oracle.  It is never consulted by an operation to
decide what the code does; it is consulted only to decide whether what the code
did was right.

% ============================================================
\section{Free Types and Constants}
% ============================================================

\begin{zed}
  [FRAME]
\end{zed}

A frame, addressed as the code addresses it: by identity, so that a frame that
vanishes and is re-created under the same identity is the same \texttt{FRAME}.
(The code pairs the frame id with its Hub; two Hubs may mint one id, and the
pair is the identity.  The model needs only that identity be a single value.)

\begin{zed}
  ZBOOL ::= ztrue | zfalse
\end{zed}

A two-valued flag as a free type, so ProB animates it as a small enumerated
set.

\begin{zed}
  STATUS ::= gone | shown | folded
\end{zed}

What the replica and the renderer know of a frame.  \texttt{gone}: the replica
holds no such frame.  \texttt{shown}: it exists and its window is expanded, so
its glyph is drawn and can be hovered.  \texttt{folded}: it exists but its
glyph is not drawn this frame --- minimized to the dock bar, put away by the
user, or its window collapsed.  The distinction between \texttt{gone} and
\texttt{folded} is the one the popup's own test draws: it asks the replica
whether the frame \emph{exists}, not whether it is drawn.

\begin{zed}
  POPUP ::= closed | opened \ldata FRAME \rdata
\end{zed}

The one info popup is either closed or open for exactly one frame.  That the
popup is open for \emph{at most one} frame is not a property to be proved: it is
what the type says, exactly as the code's single \texttt{OpenPopup | None}
says it.  Invariant~5 concerns the other half, whether that one frame still
exists.

\begin{zed}
  TARGET ::= tnone | tglyph \ldata FRAME \rdata
\end{zed}

What the pointer is over at a mouse edge: nothing of interest, or the glyph of
a frame.  A glyph can be the target only while its frame is \texttt{shown}.

\begin{zed}
  ORIGIN ::= onone | oaway | oglyph \ldata FRAME \rdata
\end{zed}

Where the button went down, for the gesture in progress: \texttt{onone} when
the button is up, \texttt{oaway} when it went down on anything other than a
glyph, \texttt{oglyph~f} when it went down on the glyph of~\texttt{f}.

\begin{zed}
  PRESS ::= nopress | gpress \ldata FRAME \cross ZBOOL \rdata
\end{zed}

The controller's press record: nothing, or the press that began on the glyph
of a frame together with whether the popup was open for that frame when it
fell.  The slot is either empty or holds a value; as a free type there is no
two-element value to constrain away.

\begin{zed}
  EDGE ::= eidle | epress | erelease
\end{zed}

What the left button did this ImGui frame.

\begin{zed}
  KEY ::= nokey | escape
\end{zed}

Whether Escape is pressed in this frame.

\begin{zed}
  COVER ::= uncovered | covered
\end{zed}

Whether a press lands on the popup window (\texttt{covered}) or outside it
(\texttt{uncovered}).

\begin{zed}
  INV ::= i1 | i2 | i3 | i4 | i5
\end{zed}

The names of the five safety invariants, so that the model can record which have
been violated in a single latch (Section~\ref{sec:state}).

% ============================================================
\section{State}
\label{sec:state}
% ============================================================

The state is flat: one schema, eleven components, and no predicate at all.  The
types carry what a predicate would otherwise have to assert, and the safety
properties (Section~\ref{sec:inv}) are deliberately \emph{absent} here, so that
a violating state remains reachable and ProB can find it.

\begin{schema}{Glyphs}
  fstat : FRAME \fun STATUS \\
  site : POPUP \\
  press : PRESS \\
  hold : ZBOOL \\
  origin : ORIGIN \\
  moving : ZBOOL \\
  drag : ZBOOL \\
  cur : ZBOOL \\
  atpress : ZBOOL \\
  shut : ZBOOL \\
  bad : \power INV
\end{schema}

The first three components are the code's.  \texttt{fstat} is the replica and
renderer's view of each frame.  \texttt{site} is the popup
(\texttt{PopupSite}), closed or open for one frame.  \texttt{press} is the
single press record of \texttt{PopupController}.

The next five describe the physical gesture and are not the code's state.
\texttt{hold} is whether the button is down.  \texttt{origin} is where it went
down.  \texttt{moving} is whether that press began a window move: it did if it
went down on a glyph that no popup covered, because then ImGui took it for a
press on the frame's title bar.  \texttt{drag} is the farthest-travelled latch:
true as soon as the pointer has strayed past the drag threshold at any moment
since the press, and never cleared until the button comes up.  \texttt{cur}
is the \emph{current} offset: true while the pointer is beyond the threshold,
false again if it returns.  The two differ exactly on a drag-away-and-return,
which is the second defect above.

\texttt{atpress} is a ghost: whether the popup was open for the pressed
frame's glyph when the press fell.  It is what Invariants~1 and~2 are judged
against.

\texttt{shut} is a shadow of ImGui's own popup stack: true when ImGui has
closed the popup but the code has not yet been told.  ImGui closes a popup on
a press outside it, and on every frame of a title-bar window move, at the
\emph{end} of the frame; the code finds out only when the next frame's
\texttt{begin\_popup} reports it shut.  For one frame, therefore, the code's
\texttt{site} says open while ImGui's popup is already gone.  That window is
exactly where the third defect lived, and it is the reason the press record
exists.  \texttt{bad} is the latch of violated invariants, written only by the
oracle and read only by the goals.

% ============================================================
\section{Three Recurring Expressions}
% ============================================================

Three expressions recur in the frame operations of Section~\ref{sec:frame}.
They are written out in each operation rather than abstracted into functions,
so that every operation can be read, and every fidelity variant edited, in one
place.

\paragraph{The surviving popup.}
\texttt{kept!} is the popup after the render pass has asked the replica
whether its frame still exists: the popup survives unless its frame is
\texttt{gone}.  This is \texttt{FrameInfoPopup.render}'s first test (``its frame
is gone'').  Replacing it by the popup itself is the buggy variant of
Section~\ref{sec:fidelity}, defect~(e).

\paragraph{Whether the popup was open for a frame.}
The press record's flag is \texttt{is\_open\_for(frame)} read at the press
edge: \texttt{ztrue} if \texttt{site~=~opened~f}, \texttt{zfalse} otherwise.

\paragraph{The audit.}
After every frame the oracle checks the two invariants that are properties of a
state rather than of a gesture.  Invariant~5 is violated if the popup is open
for a frame the replica no longer holds.  The state half of Invariant~4 is
violated if the press record names a frame the physical button did not go down
on --- which includes any record that outlives the button coming up, since then
\texttt{origin} is \texttt{onone}.  The audit adds the names of the invariants
the post-state violates to \texttt{bad}.

% ============================================================
\section{Initialization}
% ============================================================

\begin{schema}{Init}
  Glyphs'
\where
  fstat' = (\lambda f : FRAME @ gone) \\
  site' = closed \\
  press' = nopress \\
  hold' = zfalse \\
  origin' = onone \\
  moving' = zfalse \\
  drag' = zfalse \\
  cur' = zfalse \\
  atpress' = zfalse \\
  shut' = zfalse \\
  bad' = \emptyset
\end{schema}

A single initial state: no frame exists, the popup is closed, the button is up,
no press is recorded, no invariant is violated.

% ============================================================
\section{Operations Between Frames}
% ============================================================

Two kinds of event happen between ImGui frames and are visible to the code only
at the next one: the replica changing under it, and the pointer moving.

\begin{schema}{FrameChange}
  \Delta Glyphs \\
  f? : FRAME \\
  s? : STATUS
\where
  fstat~f? \neq s? \\
  fstat' = fstat \oplus \{ f? \mapsto s? \} \\
  site' = site \\
  press' = press \\
  hold' = hold \\
  origin' = origin \\
  moving' = moving \\
  drag' = drag \\
  cur' = cur \\
  atpress' = atpress \\
  shut' = shut \\
  bad' = bad
\end{schema}

\texttt{FrameChange} makes a frame appear, vanish, fold or unfold, at any
moment, including between a press and its release.  It deliberately touches
nothing else: the popup is not told, because in the code it is not.  The popup
learns that its frame has gone only when the next render pass asks.

\begin{schema}{DragOut}
  \Delta Glyphs
\where
  hold = ztrue \\
  cur = zfalse \\
  cur' = ztrue \\
  drag' = ztrue \\
  fstat' = fstat \\
  site' = site \\
  press' = press \\
  hold' = hold \\
  origin' = origin \\
  moving' = moving \\
  atpress' = atpress \\
  shut' = shut \\
  bad' = bad
\end{schema}

\begin{schema}{DragBack}
  \Delta Glyphs
\where
  hold = ztrue \\
  cur = ztrue \\
  cur' = zfalse \\
  fstat' = fstat \\
  site' = site \\
  press' = press \\
  hold' = hold \\
  origin' = origin \\
  moving' = moving \\
  drag' = drag \\
  atpress' = atpress \\
  shut' = shut \\
  bad' = bad
\end{schema}

\texttt{DragOut} moves the held pointer beyond the drag threshold; it sets
both the current offset and the farthest-travelled latch.  \texttt{DragBack}
returns the pointer to within the threshold: the current offset clears and the
latch does not.  ImGui's \texttt{mouse\_drag\_max\_distance\_sqr} is the latch;
the current offset is what the drag test compared before it was corrected
(defect~(b) of Section~\ref{sec:fidelity}).

% ============================================================
\section{The ImGui Frame}
\label{sec:frame}
% ============================================================

An ImGui frame is atomic in the model, but it has an internal order that the
code depends on, and the model keeps it:
\begin{enumerate}
  \item ImGui starts the frame.  Escape, if pressed, shuts the popup now.
  \item Each drawn glyph is read and its reading is fed to the controller.  The
    controller sees the popup as the \emph{code} holds it, which is stale by
    one frame if ImGui shut it at the end of the last: this is why the press
    record stores whether the popup was open, rather than the release asking.
  \item The popup is rendered.  The popup the code held closes if its frame is
    gone, if ImGui shut it (\texttt{shut}, from the end of the last frame or
    Escape this one), or if a close was requested in step~2.  A popup opened in
    step~2 is placed fresh, in the popup layer, and survives all of these.
  \item ImGui ends the frame.  A press that does not land on the popup shuts
    the popup that was open when it fell, and every frame of a title-bar window
    move shuts a popup that stands (assumptions A1 and A2,
    Section~\ref{sec:assume}); the code learns of it in the next frame's
    step~3.  A release shuts nothing.
  \item The controller closes out the frame: a release anywhere discards the
    press record.
\end{enumerate}
The three frame operations differ in the button edge, and so in what steps 2
and~4 do.  In every one of them \texttt{site'} is the popup after step~3,
\texttt{shut'} is what step~4 left behind, and the oracle audits the
post-state; the invariants that concern a gesture are judged in
\texttt{FrameRelease}, the frame the gesture ends on.

\texttt{FrameIdle} is a frame with no button edge: the button is up or held.
Whatever is held across it is held across a real frame in which ImGui
refocuses a moving window.

\begin{schema}{FrameIdle}
  \Delta Glyphs \\
  k? : KEY \\
  kept! : POPUP
\where
  kept! = (\IF shut = ztrue \lor (\exists f : FRAME @ \\
  \quad~ site = opened~f \land fstat~f = gone) \THEN closed \ELSE site) \\
  site' = (\IF k? = escape \THEN closed \ELSE kept!) \\
  shut' = ztrue \iff (site' \neq closed \land hold = ztrue \\
  \quad~ \land moving = ztrue) \\
  bad' = bad \cup \{ i : INV | \\
  \quad~ (i = i5 \land (\exists f : FRAME @ \\
  \quad~ \quad~ site' = opened~f \land fstat~f = gone)) \\
  \quad~ \lor (i = i4 \land (\exists f : FRAME; w : ZBOOL @ \\
  \quad~ \quad~ press = gpress(f, w) \land origin \neq oglyph~f)) \} \\
  fstat' = fstat \\
  press' = press \\
  hold' = hold \\
  origin' = origin \\
  moving' = moving \\
  drag' = drag \\
  cur' = cur \\
  atpress' = atpress
\end{schema}

\texttt{k?} is Escape.  A frame with no edge neither opens nor closes the popup
on the code's say-so; it can only be shut by ImGui (Escape, or the refocus of a
moving window, which sets \texttt{shut'} for the next frame) or by the render
pass finding the frame gone.  The refocus is the conjunct
\texttt{hold~=~ztrue~$\land$~moving~=~ztrue} of the \texttt{shut'} equation; a
popup that is already closed has nothing to shut.

\begin{schema}{FramePress}
  \Delta Glyphs \\
  t? : TARGET \\
  c? : COVER \\
  k? : KEY \\
  kept! : POPUP
\where
  hold = zfalse \\
  (\forall f : FRAME | t? = tglyph~f @ fstat~f = shown) \\
  c? = covered \implies site \neq closed \\
  hold' = ztrue \\
  drag' = zfalse \\
  cur' = zfalse \\
  (\forall f : FRAME | t? = tglyph~f @ origin' = oglyph~f) \\
  t? = tnone \implies origin' = oaway \\
  moving' = ztrue \iff (t? \neq tnone \land c? = uncovered) \\
  (\forall f : FRAME | t? = tglyph~f @ \\
  \quad~ press' = gpress(f, \IF site = opened~f \THEN ztrue \ELSE zfalse) \\
  \quad~ \land atpress' = (\IF site = opened~f \THEN ztrue \ELSE zfalse)) \\
  t? = tnone \implies press' = press \land atpress' = zfalse \\
  kept! = (\IF shut = ztrue \lor (\exists f : FRAME @ \\
  \quad~ site = opened~f \land fstat~f = gone) \THEN closed \ELSE site) \\
  site' = (\IF k? = escape \THEN closed \ELSE kept!) \\
  shut' = ztrue \iff (site' \neq closed \land c? = uncovered) \\
  bad' = bad \cup \{ i : INV | \\
  \quad~ (i = i5 \land (\exists f : FRAME @ \\
  \quad~ \quad~ site' = opened~f \land fstat~f = gone)) \\
  \quad~ \lor (i = i4 \land (\exists f : FRAME; w : ZBOOL @ \\
  \quad~ \quad~ press' = gpress(f, w) \land origin' \neq oglyph~f)) \} \\
  fstat' = fstat
\end{schema}

\texttt{t?} is what the pointer is over when the button goes down.  \texttt{c?}
is whether the press lands on the popup window itself (the popup is above the
frames, and may cover a glyph); a press that does not is a press ``outside''
and ImGui shuts the popup on it.  The press is recorded against the glyph it
fell on, with whether the popup was open for \emph{that} frame at that moment
--- read from \texttt{site}, the code's view.  A press that fell on no glyph
leaves the record as it was --- which, when the previous release was properly
handled, is empty.

The guard \texttt{c?~=~covered~$\implies$~site~$\neq$~closed} excludes the
meaningless combination of a press on a popup that is not there.  The
\texttt{moving} flag is set exactly for a press that went down on a glyph and
was not swallowed by the popup.  \texttt{shut'} is set for a press that is not
on the popup while a popup stands: ImGui shuts the popup that was open when the
press fell.  (A popup opened \emph{during} the press frame is not shut by that
press; see the assumptions of Section~\ref{sec:assume}.)

\begin{schema}{FrameRelease}
  \Delta Glyphs \\
  t? : TARGET \\
  k? : KEY \\
  acted! : ZBOOL \\
  kept! : POPUP
\where
  hold = ztrue \\
  (\forall f : FRAME | t? = tglyph~f @ fstat~f = shown) \\
  acted! = ztrue \iff (\exists f : FRAME @ t? = tglyph~f \land drag = zfalse \\
  \quad~ \land (press = gpress(f, zfalse) \\
  \quad~ \quad~ \lor (press = gpress(f, ztrue) \land site = opened~f))) \\
  (\forall f : FRAME | t? = tglyph~f \land drag = zfalse \\
  \quad~ \land press = gpress(f, zfalse) @ site' = opened~f) \\
  (\forall f : FRAME | t? = tglyph~f \land drag = zfalse \\
  \quad~ \land press = gpress(f, ztrue) \land site = opened~f @ site' = closed) \\
  kept! = (\IF shut = ztrue \lor (\exists f : FRAME @ \\
  \quad~ site = opened~f \land fstat~f = gone) \THEN closed \ELSE site) \\
  acted! = zfalse \implies \\
  \quad~ site' = (\IF k? = escape \THEN closed \ELSE kept!) \\
  hold' = zfalse \\
  origin' = onone \\
  moving' = zfalse \\
  drag' = zfalse \\
  cur' = zfalse \\
  atpress' = zfalse \\
  shut' = zfalse \\
  press' = (\IF (\exists f : FRAME; w : ZBOOL @ t? = tglyph~f \\
  \quad~ \land drag = zfalse \land press = gpress(f, w)) \\
  \quad~ \THEN nopress \ELSE press) \\
  fstat' = fstat \\
  bad' = bad \cup \{ i : INV | \\
  \quad~ (i = i5 \land (\exists f : FRAME @ \\
  \quad~ \quad~ site' = opened~f \land fstat~f = gone)) \\
  \quad~ \lor (i = i4 \land (\exists f : FRAME; w : ZBOOL @ \\
  \quad~ \quad~ press' = gpress(f, w) \land origin' \neq oglyph~f)) \\
  \quad~ \lor (i = i1 \land (\exists f : FRAME @ origin = oglyph~f \\
  \quad~ \quad~ \land t? = tglyph~f \land drag = zfalse \\
  \quad~ \quad~ \land atpress = zfalse \land site' \neq opened~f)) \\
  \quad~ \lor (i = i2 \land (\exists f : FRAME @ origin = oglyph~f \\
  \quad~ \quad~ \land t? = tglyph~f \land drag = zfalse \\
  \quad~ \quad~ \land atpress = ztrue \land site' \neq closed)) \\
  \quad~ \lor (i = i3 \land drag = ztrue \land acted! = ztrue) \\
  \quad~ \lor (i = i4 \land acted! = ztrue \\
  \quad~ \quad~ \land (\forall f : FRAME @ \\
  \quad~ \quad~ \quad~ origin \neq oglyph~f \lor t? \neq tglyph~f)) \}
\end{schema}

\texttt{t?} is what the pointer is over at the release.  The three universally
quantified lines are the controller's decision, taken in the order the code
takes it.  A release is a click on a glyph if the pointer is over it and the
farthest-travelled latch is clear (\texttt{drag~=~zfalse}); a click acts only
if the press record names \emph{that} frame.  If the record says the popup was
\emph{not} open when the press fell, the click opens it for the frame
(\texttt{open\_for}).  If it says the popup \emph{was} open, the click closes
it --- but only if the popup is still open for the frame; if ImGui shut it on
the press already there is nothing to close (the \texttt{elif} in
\texttt{\_click}).  \texttt{acted!} reports whether the controller called
\texttt{open\_for} or \texttt{request\_close} at all.

A release that does not act leaves the popup to ImGui and the render pass: it
is shut by Escape, by \texttt{shut}, or by its frame having gone, and otherwise
stays.  The release is not itself an ImGui dismissal --- ImGui dismisses on a
press --- and the refocus of a window move has ended with the button, so
\texttt{shut'} is false.  The press record is discarded on every release,
whatever happened (\texttt{end\_frame}).

The last conjunct is the oracle.  It reads the physical gesture
(\texttt{origin}, \texttt{drag}, \texttt{atpress}), which no operation uses to
decide anything, and records which invariants the outcome contradicts.

% ============================================================
\section{Invariants}
\label{sec:inv}
% ============================================================

Five safety properties must hold across all reachable states, plus
deadlock-freedom.  They are stated as predicates over \texttt{Glyphs} and its
oracle fields; how each is discharged is given in
Section~\ref{sec:verify}.  None is placed in the \texttt{Glyphs} schema: a
predicate placed there would be enforced as a guard on every successor, silently
disabling any operation that would break it and thereby masking the very defect
this model exists to catch.

\paragraph{Invariant 1 --- a click on a closed popup opens it.}
Every click on the glyph of $F$ --- a release over $F$'s glyph of a press that
began on $F$'s glyph and never strayed --- with the popup not open for $F$ when
the press fell, ends with the popup open for $F$:
\[
  i1 \notin bad.
\]
This includes the press that found the popup open for \emph{another} frame:
$F$'s popup replaces it.

\paragraph{Invariant 2 --- a click on an open popup closes it.}
Every such click with the popup open for $F$ when the press fell ends with the
popup closed, whether the click itself closed it or ImGui had already shut it on
the press:
\[
  i2 \notin bad.
\]

\paragraph{Invariant 3 --- a drag never toggles.}
A press that strayed past the drag threshold, at any moment before the release,
never causes the controller to open or close the popup.  Only ImGui's own
dismissal --- which a press outside the popup always causes --- may close it:
\[
  i3 \notin bad.
\]

\paragraph{Invariant 4 --- a release acts only on the press it came from.}
A release acts only if the button went down on that same glyph, and the press
record is a function of the physical press: it never names a frame the button
did not go down on, and it does not outlive the button coming up:
\[
  i4 \notin bad.
\]

\paragraph{Invariant 5 --- the popup is open for a frame that exists.}
At the end of every ImGui frame the popup is closed or is open for a frame the
replica still holds:
\[
  i5 \notin bad.
\]
It is judged at the end of a frame and not between frames.  A frame can vanish
at any moment, but the popup is drawn only in the render pass, and the render
pass closes it first; there is no frame in which a popup is shown for a frame
that is gone.  That the popup is open for \emph{at most one} frame is in the
type (\texttt{POPUP}).

\paragraph{Invariant 6 --- deadlock freedom.}
No reachable state leaves the machine unable to progress.  \texttt{FrameIdle}
is enabled in every state, so a frame can always be rendered.

% ============================================================
\section{Verification}
\label{sec:verify}
% ============================================================

Type-checking is a \texttt{fuzz} pass:
\begin{verbatim}
  fuzz -t docs/frame_info_popup.tex
\end{verbatim}

Invariants 1--5 are checked by asking ProB whether the negation is
\emph{reachable}.  A goal that is not found means the invariant holds on every
reachable state.  \texttt{MAX\_OPERATIONS} is raised because the frame
operations take enough inputs to exceed ProB's default of ten enabled
operations per state, and a silent truncation would make a ``not found''
meaningless:
\begin{verbatim}
  # Invariants 1-5 (each must report: No counter example found.
  #                 ALL states visited.)
  probcli docs/frame_info_popup.tex -model_check -nodead \
    -goal "i1 : bad" -p DEFAULT_SETSIZE 2 -p MAX_OPERATIONS 1000
  (and likewise i2, i3, i4, i5)
\end{verbatim}

Invariant 6 is the deadlock check over the full reachable state space:
\begin{verbatim}
  probcli docs/frame_info_popup.tex -model_check \
    -p DEFAULT_SETSIZE 2 -p MAX_OPERATIONS 1000
\end{verbatim}

Results, over two frames (the smallest carrier in which one frame's popup can
be replaced by another's) and three:
\begin{center}
\begin{tabular}{lrrl}
  carrier & states & transitions & result \\ \hline
  2 frames & 594 & 6172 & all of Invariants 1--5 hold; no deadlock \\
  3 frames & 3105 & 40582 & all of Invariants 1--5 hold; no deadlock \\
\end{tabular}
\end{center}
Every state is visited, so ``holds'' is exhaustive over the bounded carrier and
not a sample.

% ============================================================
\section{Fidelity: reproducing the defects}
\label{sec:fidelity}
% ============================================================

A model that cannot reproduce a known defect proves nothing.  Each defect below
is this specification with one rule changed, in its own file, and each then
makes one invariant's negation reachable.  Defects (a)--(d) are the four defects
of the Introduction; defect (e) is the Invariant~5 control, included because
Invariant~5 otherwise has no counterexample to show it can fail.  Each is run
breadth-first (\texttt{-bf}) so the trace is a shortest one:
\begin{verbatim}
  probcli docs/frame_info_popup_<defect>_buggy.tex -model_check -bf \
    -nodead -goal "<inv> : bad" -p DEFAULT_SETSIZE 2 -p MAX_OPERATIONS 1000
\end{verbatim}
Where the shortest trace ProB reports happens to use Escape or a press on the
popup, the trace below is the equivalent one with both excluded, which is the
user's own sequence and is also found.

\paragraph{Defect (a) --- the popup opened on the press.}
File: \texttt{frame\_info\_popup\_press\_opens\_buggy.tex}.  \texttt{FramePress}
opens the popup for the pressed glyph's frame (closing it if it was open for
that frame), and \texttt{FrameRelease} no longer acts.  Goal: \texttt{i1 : bad}.
Trace, four steps:
\begin{enumerate}
  \item \texttt{FrameChange(F, shown)}.
  \item \texttt{FramePress(uncovered, nokey, glyph of F)}: the popup opens for
    $F$.  The press is on a title bar, so the window is now moving
    (\texttt{moving = ztrue}).
  \item \texttt{FrameIdle(nokey)}: the button is held; the moving window
    refocuses and ImGui shuts the popup at the end of the frame (A2).
  \item \texttt{FrameRelease(nokey, glyph of F)}: the render pass finds the
    popup shut, and the controller no longer acts on a release.  The click on a
    closed popup has ended with it closed: \texttt{i1} is latched.
\end{enumerate}
This is the defect as described: the popup is never visible for a held click.
The trace needs the held frame, and so needs A2; without it the press-opened
popup survives and the variant is not a defect at all.

\paragraph{Defect (b) --- the drag judged by the current offset.}
File: \texttt{frame\_info\_popup\_current\_offset\_buggy.tex}.
\texttt{FrameRelease} reads \texttt{cur} where it should read \texttt{drag}.
Goal: \texttt{i3 : bad}.  Trace, five steps: \texttt{FrameChange(F, shown)};
\texttt{FramePress} on $F$'s glyph; \texttt{DragOut}; \texttt{DragBack};
\texttt{FrameRelease} on $F$'s glyph.  The pointer left and returned, so the
current offset is clear and the latch is set: the buggy release reads a click
and opens the popup, though the press was a window move.

\paragraph{Defect (c) --- the toggle decided on the release.}
File: \texttt{frame\_info\_popup\_release\_toggle\_buggy.tex}.
\texttt{FrameRelease} chooses between open and close by whether the popup is
open \emph{now}, and ignores the flag recorded at the press.  Goal:
\texttt{i2 : bad}.  Trace, six steps:
\begin{enumerate}
  \item \texttt{FrameChange(F, shown)}.
  \item \texttt{FramePress} on $F$'s glyph; \texttt{FrameRelease} on $F$'s glyph:
    the click opens the popup for $F$.
  \item \texttt{FramePress(uncovered, \ldots)} on $F$'s glyph: the popup is open
    (\texttt{atpress = ztrue}); the press is outside it, so ImGui shuts it at
    the end of this frame.
  \item \texttt{FrameIdle}: the render pass finds the popup shut and closes
    \texttt{site}.
  \item \texttt{FrameRelease} on $F$'s glyph: the buggy controller sees
    \texttt{site = closed} and opens the popup.  The second click has ended with
    the popup open.  \texttt{i2} is latched.
\end{enumerate}
The intact controller takes the decision from the press record, which still
says ``open'', finds nothing to close, and leaves the popup closed.

\paragraph{Defect (d) --- the press record not cleared by a release that is
not a click.}
File: \texttt{frame\_info\_popup\_stale\_press\_buggy.tex}.
\texttt{FrameRelease} discards the record only when it judged a click.  Goal:
\texttt{i4 : bad}.  Trace, three steps: \texttt{FrameChange(F, shown)};
\texttt{FramePress} on $F$'s glyph; \texttt{FrameRelease} over no glyph.  The
button is up and the record still names $F$: the state half of Invariant~4.  The
consequence the invariant exists to prevent follows, and ProB finds it too with
the state half switched off, so that only an \emph{action} can latch the
invariant: the same three steps, then \texttt{FramePress} over no glyph and
\texttt{FrameRelease} over $F$'s glyph.  The second release acts on the stale
record --- it opens the popup for a press that began nowhere near the glyph.

\paragraph{Defect (e) --- the popup surviving its frame.}
File: \texttt{frame\_info\_popup\_frame\_gone\_buggy.tex}.  \texttt{kept!}
ignores the replica: the popup is closed only by \texttt{shut}.  Goal:
\texttt{i5 : bad}.  Trace, five steps: \texttt{FrameChange(F, shown)}; a click
on $F$'s glyph (a press and a release) that opens the popup;
\texttt{FrameChange(F, gone)}; \texttt{FrameIdle}.  The popup is open for a
frame the replica no longer holds.

% ============================================================
\section{Assumptions About ImGui}
\label{sec:assume}
% ============================================================

The model is exact about the controller and the popup, which are the code under
test, and takes three facts about ImGui, which is not.  Each is stated with the
evidence for it and with what rests on it.

\paragraph{A1 --- a press outside the popup shuts it, at the end of that frame.}
ImGui closes the popup that is open when a press lands outside it, but only
after the frame's widgets have been drawn, so the code's \texttt{begin\_popup}
reports it shut on the \emph{following} frame.  This was measured in the
headless ImGui harness, and a regression test pins it
(\texttt{test\_imgui\_shuts\_the\_popup\_at\_the\_end\_of\_the\_press\_frame\_not\_during\_it}).
The first version of this model took the shutting to be immediate; the harness
showed otherwise, and \texttt{shut} was added.  The invariants hold under
\emph{both} timings: the immediate model had 432 states and the present one 594,
and neither has a violation.  The measurement concerns a popup that was already
open; a popup opened \emph{during} the press frame is not shut by that press
(the harness shows it surviving), and the model treats it so.

\paragraph{A2 --- a title-bar press that is moving a window shuts a standing
popup, on each frame it is held.}
This is the mechanism the code's own documentation gives for the first defect.
It was not reproduced in the headless harness: with the controller altered to open
the popup on the press, the popup stayed open through four held frames and four
dragging frames.  The model keeps it, because the first defect is stated in
these terms, but nothing about the intact controller rests on it.  The popup is
never standing during a title-bar press in the intact controller (a press outside
it shuts it by A1, and a press on it is not a title-bar press), and the model with
the rule deleted visits the same 594 states and has the same verdicts.  Defect~(a)
of Section~\ref{sec:fidelity} is therefore reproduced \emph{given} A2, and is the
one fidelity control whose real-ImGui counterpart is in doubt.

\paragraph{A3 --- Escape shuts the popup.}
The harness showed that, in this application's ImGui configuration (no keyboard
navigation), Escape does not dismiss the popup.  The model allows it anyway:
it can only add behaviours, and a conclusion drawn over a superset of the
real behaviours holds over the real ones.

% ============================================================
\section{From the Model to the Code}
% ============================================================

\texttt{press} is \texttt{PopupController.\_press}; \texttt{site} is
\texttt{PopupSite.\_open}, the popup's whole state, \texttt{None} when closed.
\texttt{FramePress} is \texttt{PopupController.observe} when the reading says
\texttt{pressed}: it stores a \texttt{GlyphPress(frame, was\_open)}.
\texttt{FrameRelease} is \texttt{observe} when the reading says \texttt{clicked}
(\texttt{\_click}) followed by \texttt{end\_frame}.  \texttt{drag} is
\texttt{GlyphReading.\_is\_undragged}'s
\texttt{mouse\_drag\_max\_distance\_sqr}; \texttt{cur} is the quantity it used to
read instead.  \texttt{kept!} and the ``dismissed by imgui'' rule are
\texttt{FrameInfoPopup.render}, and the order of steps in
Section~\ref{sec:frame} is \texttt{RenderLoop.\_render\_scene}'s: the glyphs
inside \texttt{\_render\_frames}, then \texttt{render\_popup}, whose
\texttt{end\_frame} closes the press.

\paragraph{What the model abstracts.}
A press and its release are never in the same ImGui frame; ImGui's input queue
splits a fast click.  Only one glyph is under the pointer at a time, so two
overlapping glyphs are not modelled.  A frame's owner facts, which arrive late
and are re-read each frame, change only the popup's \emph{rows}, never whether it
is open or for which frame, and are not in the state.  The drag threshold is a
single boolean.  Time passes only as frames.

\end{document}
